% CVPR 2027 submission -- SUPPLEMENTARY MATERIAL, REVIEW version (anonymous, line numbers, paper ID).
% Separate PDF (supp.pdf; due Mon Nov 23, 2026 AoE). Same kit, preamble and macros as main.tex; content: sections/supp.tex.
%
% Build: ./build.sh (or make) builds main.pdf first and then this file -> supp.pdf + render/supp-page-XX.png.
% References into the main paper: the labels of build/main.aux are imported with the prefix "M-" (package xr) and
% printed without hyperlinks (the target is in the other PDF): \maincref{sec:pcgs} -> "Sec. 3.3", \mainref{eq:smooth}
% -> "5". The main paper must therefore be compiled before the supplement.
% Rules: double-blind (no names, links to our code or machine paths); every number from a logged run;
% \tbd only for results of runs in progress (video, TFLOPs/latency).

\documentclass[10pt,twocolumn,letterpaper]{article}

%%%%%%%%% PAPER TYPE
\usepackage[review]{cvpr}        % REVIEW version, as main.tex

%% This file contains a number of tweaks that are typically applied to the main document.
%% They are not enabled by default, but can be enabled by uncommenting the relevant lines.

%%
%% Inline annotations; for predefined colors, refer to "dvipsnames" in the xcolor package:
%% https://tinyurl.com/overleaf-colors
%%
\newcommand{\red}[1]{{\color{red}#1}}
\newcommand{\todo}[1]{{\color{red}#1}}
\newcommand{\TODO}[1]{\textbf{\color{red}[TODO: #1]}}
%%
%% disable for camera ready / submission by uncommenting these lines  
%%
% \renewcommand{\TODO}[1]{}
% \renewcommand{\todo}[1]{#1}

%%
%% work harder in optimizing text layout. Typically shrinks text by 1/6 of page, enable
%% it at the very end of the writing process, when you are just above the page limit
%%
% \usepackage{microtype}

%%
%% fine-tune paragraph spacing
%%
% \renewcommand{\paragraph}[1]{\vspace{.5em}\noindent\textbf{#1.}}

%%
%% globally adjusts space between figure and caption
%%
% \setlength{\abovecaptionskip}{.5em}


%%
%% Allows "the use of \paper to refer to the project name"
%% with automatic management of space at the end of the word
%%
% \usepackage{xspace}
% \newcommand{\paper}{ProjectName\xspace}

%%
%% Commonly used math definitions
%%
% \DeclareMathOperator*{\argmin}{arg\,min}
% \DeclareMathOperator*{\argmax}{arg\,max}

%%
%% Tigthen underline
%%
% \usepackage{soul}
% \setuldepth{foobar}  % kit file, unchanged
%% macros.tex -- packages, method names and notation for the paper.
%% Loaded by main.tex after the kit's preamble.tex and BEFORE hyperref.
%%
%% Already loaded by cvpr.sty, do not load again: times, xspace, xcolor[dvipsnames], graphicx, amsmath, amssymb,
%% booktabs, natbib[numbers,sort&compress], caption, subcaption, url, etoolbox, enumitem, lineno (review mode),
%% and cleveref (\AtEndPreamble, i.e. after hyperref).  cvpr.sty also defines \eg \ie \cf \etc \vs \wrt \etal.
%% The kit's preamble.tex defines \red{..}, \todo{..} and \TODO{..} (red, in the PDF).

%% ------------------------------------------------------------------ packages
%% Engine guard: the kit loads \usepackage{times} without fontenc.  Under XeTeX/LuaTeX (TU encoding) that silently
%% falls back to Latin Modern body text (checked with tectonic 0.17), so only pdfLaTeX gives the CVPR typography.
\usepackage{iftex}
\ifPDFTeX\else
  \PackageError{macros}{Compile with pdfLaTeX (./build.sh or Overleaf's pdfLaTeX).
    XeTeX/LuaTeX replace the template's Times body font by Latin Modern}{Use pdflatex, e.g. ./build.sh}
\fi
\usepackage{microtype}  % pdfTeX: character protrusion + font expansion. Enabled from day one so that the page
                        % count seen while writing is the final one (the kit suggests it to save space).
\usepackage{bm}         % bold math, incl. Greek (\bm{\Phi})
\usepackage{multirow}
\usepackage{pifont}     % \cmark / \xmark in ablation tables
\usepackage{algorithm}      % float for the algorithm box (sections/method.tex, \label{alg:kalsea})
\usepackage[noend]{algpseudocode}  % algorithmicx syntax (\State, \If, \For, \Comment); noend saves lines
\usepackage{siunitx}    % S columns; numbers follow the text font and weight (\bfseries works in S columns)
\sisetup{mode=match, propagate-math-font=true, reset-math-version=false, reset-text-family=false,
  reset-text-series=false, reset-text-shape=false, text-family-to-math=true, text-series-to-math=true,
  group-digits=false}

%% Double-blind hygiene: do not copy file names or info dictionaries of included PDFs into the output.
\ifdefined\pdfsuppressptexinfo \pdfsuppressptexinfo=-1 \fi

%% Shared method colors (figures/src/palette.tex: cOurs, cKalSea, cSeaCache, cReference, ...), so text,
%% captions and tables can use the same colors as the figures.  Skipped silently if the file does not exist.
\InputIfFileExists{figures/src/palette.tex}{}{}

%% ------------------------------------------------------------------ title
%% Shared by main.tex and supp_main.tex (the supplement repeats the title). Candidates: TITLES.md.
\newcommand{\papertitle}{KalSea: Kalman-Corrected Spectral-Evolution-Aware Caching}

%% ------------------------------------------------------------------ method names
%% Use these macros everywhere (text, tables, captions), never the literal names.
%% Naming (owner, 2026-10-23): the full method is KalSea; the predictor alone is "KalSea w/o PCGS". Keep
%% figures/src/figstyle.py (LABEL) and figures/src/fig_qualitative.py (LABEL) in sync.
\newcommand{\ours}{KalSea\xspace}               % full method: SEA gate + Kalman residual forecasting + gap smoothing
\newcommand{\kalsea}{KalSea w/o \pcgs}          % ablation: SEA gate + Kalman residual forecasting, no gap smoothing
\newcommand{\seacache}{SeaCache\xspace}  % SEA gate + residual reuse (Chung et al., CVPR 2026)
\newcommand{\pcgs}{PCGS\xspace}          % predictor-corrector gap smoothing (the new component)
\newcommand{\sea}{SEA\xspace}            % Spectral-Evolution-Aware (filter / gate) of SeaCache
\newcommand{\nfe}{NFE\xspace}            % number of transformer (block-stack) evaluations per image
\newcommand{\flux}{FLUX.1-dev\xspace}

%% ------------------------------------------------------------------ notation
%% One symbol per quantity, everywhere.  i: sampler step (1..N); a < b: consecutive computed steps; a+1..b-1: gap.
%%   \lat_i  latent x_i            \vel_i  velocity             \sigma_i  noise level (Euler: x_{i+1} = x_i + (\sigma_{i+1}-\sigma_i) v_i)
%%   \hin_i, \hout_i  input / output of the transformer block stack \blocks (image tokens);  \cond_i  conditioning (temb)
%%   \res_i = \hout_i - \hin_i  block-stack residual;  \resf_i  forecast or smoothed residual used on skipped steps
%%   \Emb(.)  input embedding (x_embedder);  \Out(., \cond_i)  output layers (norm_out + proj_out), always run
%%   \modin_i  timestep-modulated input of the first block;  \seaop(.)  SEA filter;  \gdist_i  SEA-filtered
%%             relative L1 change;  \delta  gate threshold (skip while the accumulated \gdist stays below \delta)
%%   Kalman, per element (or per DCT band k):  state \kstate_i = [m_i, \nu_i]^T (value, derivative), warped time s,
%%             \ds_i step, \rho = exp(-\ds/\tau), \kA transition, \kP covariance, \kQ = q \kappa [[ds^3/3, ds^2/2],
%%             [ds^2/2, ds]] process noise (\kappa adapted online), R measurement noise, z_b observation,
%%             e_b = z_b - \prior{m}_b innovation, \kPhi(b,i) transition from i to b, g_i smoother gain
%%   \prior{m}_i  predicted (prior);  \post{m}_i  filtered (posterior);  \smooth{m}_i  RTS-smoothed
\newcommand{\lat}{\mathbf{x}}
\newcommand{\vel}{\mathbf{v}}
\newcommand{\hin}{\mathbf{h}^{\mathrm{in}}}
\newcommand{\hout}{\mathbf{h}^{\mathrm{out}}}
\newcommand{\res}{\mathbf{r}}
\newcommand{\resf}{\tilde{\mathbf{r}}}
\newcommand{\cond}{\mathbf{c}}
\newcommand{\blocks}{\mathcal{F}}
\DeclareMathOperator{\Emb}{Emb}
\DeclareMathOperator{\Out}{Out}
\DeclareMathOperator{\seaop}{SEA}
\DeclareMathOperator{\DCT}{DCT}
\newcommand{\modin}{\mathbf{I}}
\newcommand{\gdist}{d}
\newcommand{\kstate}{\mathbf{s}}
\newcommand{\kA}{\mathbf{A}}
\newcommand{\kP}{\mathbf{P}}
\newcommand{\kQ}{\mathbf{Q}}
\newcommand{\kPhi}{\bm{\Phi}}
\newcommand{\ds}{\Delta s}
\newcommand{\prior}[1]{#1^{-}}
\newcommand{\post}[1]{#1^{+}}
\newcommand{\smooth}[1]{#1^{\mathrm{s}}}
\newcommand{\E}{\mathbb{E}}
\DeclareMathOperator*{\argmin}{arg\,min}
\DeclareMathOperator*{\argmax}{arg\,max}

%% ------------------------------------------------------------------ tables
\newcommand{\best}[1]{\textbf{#1}}          % best value in a column
\newcommand{\second}[1]{\underline{#1}}     % second best
\newcommand{\nodata}{--}                    % intentionally blank cell (TFLOPs / latency until measured)
\newcommand{\tbd}{\textcolor{red}{TBD}\xspace} % result not available yet (video); must be gone before submission
\newcommand{\cmark}{\ding{51}}
\newcommand{\xmark}{\ding{55}}

%% ------------------------------------------------------------------ placeholders for parallel work
%% \maybeinput{tables/foo}  inputs tables/foo.tex if it exists, else draws a red box.
%% \maybegraphics[opts]{figures/foo.pdf}  includes the figure if it exists, else draws a red box of the same width.
%% Missing files are logged as "MISSING-FILE: <name>"; build.sh / tools/pdfcheck.py report them.
\newcommand{\missingbox}[2]{%
  \typeout{MISSING-FILE: #2}%
  \fbox{\parbox[c][1.5cm][c]{\dimexpr#1-2\fboxsep-2\fboxrule\relax}{\centering\footnotesize\color{red}%
    missing: \texttt{\detokenize{#2}}}}}
\newcommand{\maybeinput}[1]{\IfFileExists{#1.tex}{\input{#1.tex}}{\IfFileExists{#1}{\input{#1}}{%
  \missingbox{\linewidth}{#1}}}}
\newcommand{\maybegraphics}[2][width=\linewidth]{\IfFileExists{#2}{\includegraphics[#1]{#2}}{%
  \missingbox{\linewidth}{#2}}}
    % our packages, title, method names, notation

\usepackage{placeins}            % \FloatBarrier: all floats before the references
\usepackage{xr}                  % labels of the main paper (before hyperref)
\externaldocument[M-]{build/main}

\definecolor{cvprblue}{rgb}{0.21,0.49,0.74}
\usepackage[pagebackref,breaklinks,colorlinks,allcolors=cvprblue]{hyperref}

\newcommand{\maincref}[1]{\cref*{M-#1}}  % "Sec. 3.3", "Tab. 1", "Fig. 3" of the main paper, no hyperlink
\newcommand{\mainref}[1]{\ref*{M-#1}}    % bare number of a main-paper label (e.g. an equation)
\newcommand{\maineqref}[1]{(\ref*{M-#1})}

% Supplementary numbering: sections A, B, ...; tables, figures and equations S1, S2, ... (main-paper items keep their
% plain numbers, so "Tab. 1" is the main paper and "Tab. S1" this document).
\renewcommand{\thesection}{\Alph{section}}
\renewcommand{\thetable}{S\arabic{table}}
\renewcommand{\thefigure}{S\arabic{figure}}
\renewcommand{\theequation}{S\arabic{equation}}

% Float placement: the supplement is mostly tables; allow more and larger floats per page than the LaTeX defaults
% so that pages are not left half empty (the main paper keeps the defaults).
\renewcommand{\topfraction}{0.9}
\renewcommand{\dbltopfraction}{0.9}
\renewcommand{\bottomfraction}{0.6}
\renewcommand{\textfraction}{0.08}
\renewcommand{\floatpagefraction}{0.75}
\renewcommand{\dblfloatpagefraction}{0.75}
\setcounter{topnumber}{4}
\setcounter{dbltopnumber}{4}
\setcounter{bottomnumber}{2}
\setcounter{totalnumber}{6}
\makeatletter  % float-only pages: stack the floats from the top instead of spreading them over the page
\setlength{\@fptop}{0pt}\setlength{\@fpsep}{14pt plus 2pt}\setlength{\@fpbot}{0pt plus 1fil}
\setlength{\@dblfptop}{0pt}\setlength{\@dblfpsep}{14pt plus 2pt}\setlength{\@dblfpbot}{0pt plus 1fil}
\makeatother

\raggedbottom  % the supplement has many float pages; do not stretch the text columns to the full height

% Auto-generated by figures/src/make_supp_tables.py from logged evaluation files -- edit the generator, not this file.
% statistics quoted in sections/supp.tex (paired over prompts, bootstrap as in the harness)
\newcommand{\suppKalRerunA}{+0.03\,dB [$-$0.003, +0.07]}
\newcommand{\suppKalRerunSameA}{200}
\newcommand{\suppKalRerunB}{+0.002\,dB [$-$0.01, +0.01]}
\newcommand{\suppKalRerunSameB}{200}
\newcommand{\suppPcgsMedianA}{0.57}
\newcommand{\suppPcgsNegA}{39}
\newcommand{\suppPcgsWorstA}{4.75}
\newcommand{\suppPcgsMedianB}{0.44}
\newcommand{\suppPcgsNegB}{44}
\newcommand{\suppPcgsWorstB}{2.31}
  % statistics quoted in the text (generated by figures/src/make_supp_tables.py)

%%%%%%%%% PAPER ID (same as main.tex)
\def\paperID{*****} % OpenReview paper ID, available after registration (Nov 10, 2026 AoE); keep equal to main.tex
\def\confName{CVPR}
\def\confYear{2027}

\title{\papertitle}
\author{Anonymous}

\begin{document}
\maketitlesupplementary   % kit command: title + "Supplementary Material" (no author block)
%% sections/supp.tex -- OUTLINE and first draft of the supplementary material.
%% The supplement is a separate, anonymized PDF (due Mon Nov 23, 2026 AoE); main.tex does NOT input this file.
%% Intended use: a supp.tex driver with the same preamble/macros that inputs this file (references to the main
%% paper are written as "main paper, Sec. X" because labels of the main PDF are not available here).
%% Content marked \TODO is planned but not written yet; every number must come from a logged run.

\section{Derivation of the gap smoother}
\label{supp:smoother}

\paragraph{Setting.}
Consider one Kalman filter of the predictor (one residual element, or one DCT band whose coefficients share the
covariance). Let $a<b$ be consecutive computed steps, $\mathcal Y_a$ the observations up to step $a$, and $a<i<b$ a
skipped step. No observation falls into the gap, so the predicted moments are propagated without update:
$\prior{\kstate}_{i}=\kA_i\prior{\kstate}_{i-1}$ and $\prior{\kP}_i=\kA_i\prior{\kP}_{i-1}\kA_i^\top+\kQ_i$ for
$a+1<i\le b$, starting from $\prior{\kstate}_{a+1}=\kA_{a+1}\post{\kstate}_a$. Define the transition
$\kPhi(b,i)=\kA_b\kA_{b-1}\cdots\kA_{i+1}$, with $\kPhi(b,b)=\mathbf I$, and the innovation variance
$S_b=\mathbf h^\top\prior{\kP}_b\mathbf h+R$.

\paragraph{Proposition.}
For $a<i\le b$,
\begin{equation}
  \E\!\left[\kstate_i\mid\mathcal Y_a,z_b\right]=\prior{\kstate}_i+\prior{\kP}_i\kPhi(b,i)^\top\mathbf h\,S_b^{-1}e_b,
  \qquad e_b=z_b-\mathbf h^\top\prior{\kstate}_b,
  \label{eq:supp-smooth}
\end{equation}
and the first component gives the gain $g_i=[\prior{\kP}_i\kPhi(b,i)^\top]_{00}/S_b$ of the main paper.

\paragraph{Proof by conditioning.}
Iterating the state equation, $\kstate_b=\kPhi(b,i)\kstate_i+\mathbf w_{i\to b}$, where $\mathbf w_{i\to b}$ is a sum of
process-noise terms of steps $i+1,\dots,b$; it has zero mean and is independent of $\kstate_i$ and of $\mathcal Y_a$.
With $z_b=\mathbf h^\top\kstate_b+\eta_b$ and $\eta_b\sim\mathcal N(0,R)$ independent of everything else,
$\mathrm{Cov}(\kstate_i,z_b\mid\mathcal Y_a)=\prior{\kP}_i\kPhi(b,i)^\top\mathbf h$ and
$\mathrm{Var}(z_b\mid\mathcal Y_a)=S_b$. Conditioning the jointly Gaussian pair $(\kstate_i,z_b)$ given $\mathcal Y_a$
on $z_b$ gives \eqref{eq:supp-smooth}.

\paragraph{Equivalence with the RTS smoother.}
The RTS recursion~\cite{rauch1965maximum,sarkka2013bayesian} over $[a,b]$, using the observations up to $b$, reads
$\smooth{\kstate}_i=\post{\kstate}_i+\mathbf C_i(\smooth{\kstate}_{i+1}-\prior{\kstate}_{i+1})$ with
$\mathbf C_i=\post{\kP}_i\kA_{i+1}^\top(\prior{\kP}_{i+1})^{-1}$, initialized with the filtered estimate
$\smooth{\kstate}_b=\post{\kstate}_b=\prior{\kstate}_b+\prior{\kP}_b\mathbf h S_b^{-1}e_b$. Inside the gap,
$\post{\kstate}_i=\prior{\kstate}_i$ and $\post{\kP}_i=\prior{\kP}_i$. We show by backward induction that
$\smooth{\kstate}_i-\prior{\kstate}_i=\prior{\kP}_i\kPhi(b,i)^\top\mathbf h S_b^{-1}e_b$. The case $i=b$ holds because
$\kPhi(b,b)=\mathbf I$. If it holds for $i+1$, then
\begin{equation*}
  \smooth{\kstate}_i-\prior{\kstate}_i
  =\prior{\kP}_i\kA_{i+1}^\top(\prior{\kP}_{i+1})^{-1}\prior{\kP}_{i+1}\kPhi(b,i+1)^\top\mathbf h S_b^{-1}e_b
  =\prior{\kP}_i\big(\kPhi(b,i+1)\kA_{i+1}\big)^\top\mathbf h S_b^{-1}e_b,
\end{equation*}
and $\kPhi(b,i+1)\kA_{i+1}=\kPhi(b,i)$. \hfill$\square$

\paragraph{Consequences.}
\begin{itemize}
  \item At the newest observation the smoothed and filtered estimates coincide, so \pcgs leaves the predictor for the
  next gap unchanged (main paper, Remarks of the \pcgs section).
  \item \emph{Random walk.} For the model without trend ($\kA=1$, $\kQ=q\ds$), an exact observation at $a$
  ($\post{\kP}_a=0$) and $R\to0$, the gain is $g_i=(s_i-s_a)/(s_b-s_a)$, \ie, linear interpolation in the time
  coordinate. With the reuse forecast $\hat\res_i=\res_a$ this is the linear interpolation between $\res_a$ and
  $\res_b$ used for \seacache + \pcgs in the ablations.
  \item \emph{DCT variant.} Each band $\ell$ has its own $\kA$, $\kQ$, $\kP$ and hence its own gains $g_{i,\ell}$. Band
  normalization cancels in \eqref{eq:supp-smooth}, so the correction in residual space is
  $\resf_i=\hat\res_i+\mathcal D^{-1}\big(\mathbf G_i\odot\mathcal D(\res_b-\hat\res_b)\big)$ with the per-channel
  orthonormal 2D DCT $\mathcal D$ and the gain map $\mathbf G_i$ (value $g_{i,\ell}$ on band $\ell$).
  \item \emph{Gain shape.} \TODO{Figure: $g_i$ per band for the modal 13-NFE schedule of \flux
  (0,1,2,4,6,9,13,18,25,33,41,47,49). Preliminary computation with the \flux calibration and $\kappa=1$ (no online
  adaptation): within the seven-step gaps the gain rises monotonically from 0.04--0.19 at the first skipped step to
  0.74--0.88 at the last one (high-frequency bands more convex); recompute with the logged $\kappa$.}
  \item \emph{First-order optimality.} \TODO{Write out the linearized error propagation
  $\lat_b-\lat^{\mathrm{ref}}_b\approx\sum_i\mathbf M_{b,i}(\sigma_{i+1}-\sigma_i)\mathbf J_i(\hat\res_i-\res_i)$ and the
  conditional-mean argument of the main paper.}
\end{itemize}

\section{Forecaster details}
\label{supp:forecaster}
\begin{itemize}
  \item DCT bands: orthonormal DCT-II per channel on the $64{\times}64$ token grid; seven radial bands with normalized
  radius $r=\sqrt{(u/H)^2+(v/W)^2}/\sqrt2$ and edges $\{0,\tfrac1{64},\tfrac1{32},\tfrac1{16},\tfrac18,\tfrac14,
  \tfrac12,1\}$; band normalization by the calibrated residual power.
  \item Time coordinate (following SpecKCast~\cite{rajabi2026speckcast}): $\lambda_i=2\log((1-\sigma_i)/\sigma_i)$
  (clamped), $t_\ell(i)=\mathrm{sigmoid}((\lambda_i+\log p_\ell)/T_w)$ with $T_w=2$, rescaled to $[-1,1]$; $p_\ell$ is
  the calibrated power of the clean-sample estimate in band $\ell$. Per-element variant: one band with the
  count-weighted mean power.
  \item Calibration (16 COCO captions, disjoint from the 64 development and the 200 DrawBench prompts): residual band
  powers, $p_\ell$ from the last 25 steps, process-noise intensity $q_\ell$ as the median over steps of the mean squared
  change of the finite-difference trend per unit warped time, initial trend variance $2q_\ell$.
  \item Initialization at the first computed step: $m=z_0$, $\nu=0$, $\kP=\mathrm{diag}(R,2q)$.
  \item Process-noise adaptation: $\gamma_b=\mathrm{clip}_{[0.1,10]}(\overline{e_b^2}/S_b)$ with the band-mean squared
  innovation, $\kappa\gets\mathrm{clip}_{[0.1,10]}(\kappa\gamma_b^{0.2})$.
  \item Calibration-free variant: unit band powers, $q=7400$, initial trend variance $15000$, uniform step time.
  \TODO{state which variant is used for video.}
\end{itemize}

\section{Implementation of \pcgs}
\label{supp:impl}
\begin{itemize}
  \item Bookkeeping per skipped step: latent $\lat_i$ (only $\lat_{a+1}$ is needed for the exact replay), conditioning
  $\cond_i$, forecast $\hat\res_i$, and the prior covariance of every band; memory cost (\flux): \TODO{MiB}.
  \item Exact replay: same tensors and arithmetic as the diffusers flow-matching Euler step (fp32 accumulation of the
  bf16 product, cast to bf16); with the identity smoother ($\resf_i=\hat\res_i$) the replay reproduces the uncorrected
  run bitwise (unit test).
  % provenance: research_v2/ablations/ABLATIONS_tests.md, test U3; T1 for the next item.
  \item The implemented gains match brute-force joint-Gaussian conditioning of the damped-trend model to a relative
  error below $10^{-15}$ for both parametrizations (unit test T1, float64, synthetic residuals, \flux calibration).
  \item Continuation at $b$ and the gate refresh (first-block adaptive normalization and \sea filter); the last step is
  corrected but needs no refresh.
  \item NFE accounting: forward pre-hooks on all 19+38 blocks, cross-checked against the forward pass (zero mismatches,
  zero partial block stacks in every run).
  \item HunyuanVideo port: Euler sampler, single guidance-distilled branch; the official gate binds the scheduler that
  exists at load time, so its \sea filter reads $\sigma=1-i/1000$ instead of the sampling schedule. We reproduce this
  behavior for all gated methods (\seacache's thresholds were tuned with it) and use the live schedule for the
  re-integration. \TODO{confirm with the final port.}
  \item Wan2.1 port: UniPC multistep solver with classifier-free guidance; per-branch filters; solver-state snapshot at
  the start of every gap and replay of the gap with the corrected outputs. \TODO{validate; report the bitwise test.}
\end{itemize}

\section{Evaluation protocol}
\label{supp:protocol}
\begin{itemize}
  \item DrawBench order and seeds (seed = prompt index); reference = uncached 50-step output of the same code path.
  \item Budget mapping of the official \sea gate on \flux: $\delta=0.2/0.3/0.37/0.5/0.6/0.8/1.0\to
  27/20.8/18.1/15/13/11/9$ \nfe. \seacache's Table~1 labels its $\approx$50\% row $\delta=0.3$ but its TFLOPs imply
  18.45 full steps, which matches our $\delta\approx0.37$; its $\delta=0.6$ row (12.99 steps) matches ours (13.00).
  We therefore match budgets by measured \nfe.
  \item SSIM variant (scikit-image, $7{\times}7$ uniform window) vs.\ the Gaussian SSIM of TeaCache's evaluation code;
  LPIPS (AlexNet, inputs in $[-1,1]$); PSNR on 8-bit RGB.
  \item Video: 944 unique VBench prompts (two duplicates in VBench's list of 946); frame counts (HunyuanVideo 65, Wan2.1
  81, the latter matching \seacache's reported TFLOPs); metrics on decoded MP4 frames with TeaCache's metric code;
  \TODO{seeds}.
  \item Bootstrap: paired over prompts, 10{,}000 resamples, percentile intervals, identical resampling indices for all
  pairs.
\end{itemize}

\section{Additional results}
\label{supp:results}
\begin{itemize}
  \item Full budget sweep with confidence intervals for all fillers, including linear extrapolation
  (table view of the main-paper trade-off figure plus a linear-extrapolation column).
  % source: tables/tab_flux_sweep.tex (generated); add the respect_<d>_w0 runs as 'linear extrapolation'.
  \item Open-loop forecast error inside gaps on the development set (six reference trajectories): relative velocity
  error per skipped step for the residual, clean-image, noise and velocity targets. \TODO{figure.}
  \item Correction size $\|\lat'_b-\lat_b\|/\|\lat_b\|$ versus gap length and position, all budgets. \TODO{figure.}
  \item Development-set (COCO-64) results for all configurations, with the halves used for development and
  confirmation. \TODO{development-set table.}
  % source: research_v2/ablations/ABLATIONS_devset.md
  \item Full ablation table, including the continuation variants.
  % source: tables/tab_ablation_full.tex (generated by research_v2/ablations/make_tables.py)
  \item Video: per-dimension VBench scores and per-budget tables. \TODO{after the runs.}
  \item TFLOPs and latency for all methods. \TODO{after the controlled measurement.}
\end{itemize}

\section{Negative and neutral results}
\label{supp:negative}
\begin{itemize}
  \item Forecast target: forecasting the clean-image, noise or velocity prediction reconstructed from the current
  latent, instead of the block-stack residual, is not better on the development set (all differences within about
  $\pm$0.13\,dB, intervals include zero).
  \item Allocation: an innovation-driven controller that replaces the \sea gate. \TODO{Owner decision: the interim
  development-set data show a large schedule (allocation) effect that a fixed schedule reproduces, and the predictive
  variant fails; this is not evidence that error control is worse than the \sea gate. Decide whether to report it.}
\end{itemize}

\section{More qualitative results}
\label{supp:qualitative}
\begin{itemize}
  \item Qualitative results at 13 and 20.8 \nfe with the selection rule of the main paper and more quantiles.
  % figures/fig_qualitative_supp_nfe13.pdf, figures/fig_qualitative_supp_nfe21.pdf
  \item Video frames. \TODO{after the runs.}
\end{itemize}

\FloatBarrier\twocolumn  % \twocolumn: fresh page state after the float pages (avoids a stretched first reference column)
{
    \small
    \bibliographystyle{ieeenat_fullname}
    \bibliography{refs}
}
\end{document}
